\section{Introduction}
Fractional vegetation cover (FVC) supports ecosystem and land-degradation assessment in arid and alpine settings \citep{chen2016uav,li2019plateau,han2023}. NDVI \citep{tucker1979} is widely used in calibrated FVC mappings and in the dimidiate pixel model (DPM), a two-endmember formulation \citep{carlson1997,gutman1998,jiapaer2011,yan2022dpm}. DPM transparency is offset by sensitivity to background, endmembers, shadow, and departures from two-component mixing \citep{huete1988,montandon2008,mu2024}; calibrated multi-source retrievals instead require matched references and transfer assessment \citep{wang2018fvc,mu2021multisource,song2022fvc}.

Copernicus FCOVER V2 RT6 is a 300 m, 10-day operational green-cover product with documented retrieval and validation history \citep{baret2013,camacho2013,liu2019,fuster2020,copernicus2025validation,copernicus2026atbd,copernicus2026pum}. Sentinel-2, Landsat, and MODIS differ in support, spectral response, timing, geometry, and QA \citep{drusch2012,roy2014,vermote2016}; residual NDVI differences remain after product-level correction \citep{claverie2018,zhang2018ndvi,shang2019,trevisiol2024}. Because observation support affects remotely sensed values and variance \citep{woodcock1987,atkinson2000}, direct comparison of native pixels with 300 m FCOVER can mix retrieval disagreement with grid mismatch.

Existing work generally treats FVC/DPM formulation, FCOVER validation, sensor harmonization, and spatial or chronological validation separately \citep{roberts2017,meyer2018,ploton2020,tashman2000,bergmeir2012}. The present study integrates common-target-grid harmonization, geographic replication, block-grouped diagnostics, and chronological transfer to test whether a simple NDVI--FCOVER relation is stable across those dimensions.

We evaluate (RQ1) common-target-grid FCOVER-reference agreement and the empirical-quantile DPM comparator across four AOIs; (RQ2) variation in error, history preference, and OLS coefficients; and (RQ3) whether longer histories are uniformly associated with lower 2024/2025 forward error. The design is an evaluation contribution, not a new retrieval algorithm: it neither identifies a causal effect of added years nor establishes field-level FVC accuracy.
